\documentclass[10pt,a4paper]{article}

\usepackage[margin=1.55cm]{geometry}
\usepackage{amsmath,amssymb}
\usepackage{array,booktabs,longtable}
\usepackage{enumitem}
\usepackage{xcolor}
\usepackage[hidelinks]{hyperref}
\usepackage{microtype}
\setlength{\parindent}{0pt}
\setlength{\parskip}{0.52em}
\setlength{\LTpre}{0.25em}
\setlength{\LTpost}{0.75em}
\renewcommand{\arraystretch}{1.12}

\newcommand{\ttbar}{t\bar t}
\newcommand{\ttW}{t\bar t W}
\newcommand{\pt}{p_{\mathrm T}}
\newcommand{\met}{p_{\mathrm T}^{\mathrm{miss}}}
\newcommand{\ddphi}{\Delta\phi}
\newcommand{\dr}{\Delta R}
\newcommand{\Wp}{W^+}
\newcommand{\Wm}{W^-}
\newcommand{\code}[1]{\texttt{\detokenize{#1}}}
\newcolumntype{L}[1]{>{\raggedright\arraybackslash}p{#1}}

\title{Comprehensive fixed-order analysis of \(t\bar tW^\pm\)\\
in the trilepton channel\\
\large Observable definitions and implementation reference}
\author{fNLO fixed-order template}
\date{30 August 2026}

\begin{document}
\maketitle

\begin{abstract}
This document specifies the 286-histogram HwU analysis
\code{analysis_HwU_pp_ttxw_trilepton.o} for
\[
 pp\longrightarrow t\bar tW^\pm+X
 \longrightarrow b\,\bar b\;e\,e\,\mu\;3\nu+X .
\]
The no-cut truth layer uses the electron and positron to analyse the top and
antitop decays and the muon to analyse the associated \(W\). The
experiment-facing layer is formulated only in terms of the observable
two-electron/one-muon final state, jets, and total missing transverse
momentum. It includes both \(e^+e^-\mu^\pm\) and
\(e^\pm e^\pm\mu^\mp\) charge patterns and does not assume which parent W
produced a visible lepton.
It includes the six distributions unfolded in the ATLAS differential
measurement as well as extensive production, decay, polarization,
spin-correlation, event-activity, transverse, and neutrino-blind
reconstruction observables. One IR-safe measurement function is used for
additive and multiplicative NLO decay-chain events and for every real-event
counterevent.
\end{abstract}

\tableofcontents

\section{Scope and event interpretation}

The flavour-specific theory interpretation is
\begin{align*}
 \ttbar W^+&\to b\,e^+\nu_e\;\bar b\,e^-\bar\nu_e\;\mu^+\nu_\mu,\\
\ttbar W^-&\to b\,e^+\nu_e\;\bar b\,e^-\bar\nu_e\;\mu^-\bar\nu_\mu .
\end{align*}
The event-level interface is more general: it requires exactly two
status-one electron-flavour charged leptons, one muon, three light
neutrinos, one \(b\), and one \(\bar b\). The charged-lepton sum must be
\(q_{3\ell}=\pm1\), with two leptons of charge \(q_{3\ell}\) and one of the
opposite charge. Thus all charge/origin assignments of the \(2e1\mu\)
channel are accepted. Missing or duplicate required objects stop the run
with a diagnostic.

At truth level the decay assignment is intentionally flavour-specific:
\(e^+\nu_e\) reconstructs the top-side \(\Wp\), \(e^-\bar\nu_e\) the
antitop-side \(\Wm\), and the unique muon and muon neutrino the associated
\(W^\pm\). These truth histograms are filled only when that exact PDG
content is present. In particular, a \(\mu^+\) (PDG \(-13\)) is paired with
\(\nu_\mu\) (PDG \(14\)), and a \(\mu^-\) (PDG \(13\)) with
\(\bar\nu_\mu\) (PDG \(-14\)). Identification uses PDG codes rather than
positions, so MadGraph process ordering is irrelevant.

The realistic layer does not consult ancestry or individual neutrino
momenta. It sees two electron-flavour leptons, one muon, their charges,
jets, and \(\vec p_{\mathrm T}^{\,\mathrm{miss}}\). Consequently it remains
defined when either a top decay or the associated W supplies the muon.

The analysis has three nested levels.
\begin{enumerate}[leftmargin=1.8em]
\item \textbf{Inclusive truth} has no acceptance cuts and uses exact neutrino
  momenta. It provides production and decay kinematics, perfect top/W
  reconstruction, polarization, and spin-density moments for the
  flavour-specific decay assignment above.
\item \textbf{Fiducial visible} uses only charged leptons, anti-\(k_T\) jets,
  ideal partonic b tags, and total missing transverse momentum. It requires
  at least one accepted b jet and is the level closest to a particle-level
  measurement.
\item \textbf{Neutrino-blind reconstructed} additionally requires two
  distinct accepted b jets. A deterministic three-neutrino mass fit
  reconstructs all W bosons, both top quarks, and the full \(\ttW\) system
  without reading a truth neutrino.
\end{enumerate}

All active HwU weights are filled. Histogram integrals therefore retain the
fixed-order cross-section normalization for the central, scale, PDF, and
other configured weights; the module does not normalize shapes.

\section{Kinematic conventions}

Four-vectors are \(p=(E,p_x,p_y,p_z)\), with
\begin{align}
 \pt(p)&=\sqrt{p_x^2+p_y^2},&
 y(p)&=\frac12\ln\frac{E+p_z}{E-p_z},\\
 \eta(p)&=\frac12\ln\frac{|\vec p|+p_z}{|\vec p|-p_z},&
 m(p)&=\sqrt{\max(0,p^2)} .
\end{align}
Azimuthal differences are wrapped into \([-\pi,\pi]\); unsigned
\(\ddphi\) lies in \([0,\pi]\), and
\(\dr=\sqrt{(\Delta\eta)^2+(\ddphi)^2}\). Leading, subleading, and third
refer to descending charged-lepton \(\pt\).

The associated-W charge is inferred without ancestry as
\(q_W=q_{3\ell}=\sum_{i=1}^3q_{\ell_i}=+1\) or \(-1\). The unique
same-sign pair comprises the two leptons with charge \(q_W\); they are
ordered by \(\pt\). The remaining lepton has charge \(-q_W\). A
``leading opposite-sign pair'' combines this remaining lepton with the
leading same-sign lepton. This convention is unambiguous also for
\(e^\pm e^\pm\mu^\mp\), where two opposite-sign \(e\mu\) pairs exist. The
dielectron pair is OSSF only when its charges are opposite. Multiplying an
observable by \(q_W\) permits direct charge-channel comparisons.

The missing transverse vector is computed from visible status-one particles,
\begin{equation}
 \vec p_{\mathrm T}^{\,\mathrm{miss}}
 =-\sum_{i\notin\{\nu_e,\nu_\mu,\nu_\tau\}}\vec p_{\mathrm T,i}.
\end{equation}
For a consistent fixed-order point it equals the vector sum of all three
neutrinos. Their difference is retained as a closure diagnostic, never as
a selection.

For visible \(v\) and massless transverse momentum \(\vec q_{\mathrm T}\),
\begin{align}
 m_{\mathrm T}^2(v,\vec q_{\mathrm T})
 &=m_v^2+2\left(E_{\mathrm T,v}|\vec q_{\mathrm T}|
 -\vec p_{\mathrm T,v}\cdot\vec q_{\mathrm T}\right),\\
 E_{\mathrm T,v}&=\sqrt{m_v^2+\pt^2(v)} .
\end{align}
The trilepton cluster mass uses \(v=e_1+e_2+\mu\) and treats all missing
momentum as one massless transverse object. The visible \(m_{\mathrm T2}^{ee}\)
is
\begin{equation}
 \min_{\vec q_{1\mathrm T}+\vec q_{2\mathrm T}=\vec p_{\mathrm T}^{\rm miss}}
 \max\{m_{\mathrm T}(e_1,\vec q_{1\mathrm T}),
       m_{\mathrm T}(e_2,\vec q_{2\mathrm T})\}.
\end{equation}
It is measurable but is not expected to have the clean dileptonic-\(\ttbar\)
endpoint because \(\vec p_{\mathrm T}^{\rm miss}\) also contains the
associated-W neutrino.

\section{Truth resonances and decay radiation}

The three exact W candidates are
\[
 p_{W_t^+}=p_{e^+}+p_{\nu_e},\qquad
 p_{W_{\bar t}^-}=p_{e^-}+p_{\bar\nu_e},\qquad
 p_{W_{\rm assoc}}=p_\mu+p_{\nu_\mu^{(-)}} .
\]
At Born level \(p_t=p_{W_t^+}+p_b\) and
\(p_{\bar t}=p_{W_{\bar t}^-}+p_{\bar b}\). For a real-emission event,
each non-bottom status-one light quark or gluon is assigned to production,
the top decay, or the antitop decay. The selected member of the
\(3^{N_{\rm rad}}\) possibilities minimizes
\begin{equation}
 \chi^2_{\rm truth}=
 \left(\frac{m_{t,\rm cand}-m_t^{\rm ref}}{5~\mathrm{GeV}}\right)^2+
 \left(\frac{m_{\bar t,\rm cand}-m_t^{\rm ref}}{5~\mathrm{GeV}}\right)^2 .
\end{equation}
The associated W is not part of this QCD-radiation assignment. Reference
masses are read from the model. This construction is independent of event
record ordering and approaches the same measurement in every soft or
collinear limit, apart from measure-zero assignment boundaries.

The top-side W helicity angles boost the charged lepton to its parent-W rest
frame and use the W flight direction in the corresponding top rest frame.
For the truth associated W, the helicity axis is its flight direction in the
\(\ttW\) rest frame and the analyser is the muon in its W rest frame. The
reconstructed version uses whichever same-sign lepton the mass fit assigns
to the associated W. Raw and \(q_W\)-signed angles are both plotted.

\section{Jets and event activity}

All status-one quarks with \(|\mathrm{PDG}|\leq5\) and gluons are clustered
with FastJet using anti-\(k_T\), \(R=0.4\), and four-vector recombination.
Counted jets satisfy
\[
 \pt(j)\geq25~\mathrm{GeV},\qquad |\eta(j)|<2.5,
\]
matching the central particle-level jet definition of the ATLAS
measurement. A jet is ideally b-tagged if it contains the unique final-state
\(b\) or \(\bar b\) parton. If both bottom partons merge, inclusive and
one-b fiducial plots remain defined, while the two-b reconstruction tier
fails. Additional jets are accepted jets other than the two b-jet indices.

The activity variables are
\begin{align}
 H_{\mathrm T}^{j}&=\sum_{\rm accepted~jets}\pt(j),&
 H_{\mathrm T}^{\ell}&=\pt(e_1)+\pt(e_2)+\pt(\mu),\\
 H_{\mathrm T}^{\rm extra}&=\sum_{j\notin b,\bar b}\pt(j),&
 H_{\mathrm T}^{\rm vis}&=H_{\mathrm T}^{j}+H_{\mathrm T}^{\ell},\\
 S_{\mathrm T}&=H_{\mathrm T}^{\rm vis}+\met,&
 p_{\rm vis}&=p_{e_1}+p_{e_2}+p_\mu+p_{b{\rm j}}+p_{\bar b{\rm j}} .
\end{align}
The analysis explicitly includes the six ATLAS unfolded variables:
\(N_{\rm jets}\), \(H_{\mathrm T}^{j}\), \(H_{\mathrm T}^{\ell}\),
\(\dr_{\ell b,\rm lead}\), \(|\ddphi_{\ell\ell,\rm SS}|\), and
\(|\Delta\eta_{\ell\ell,\rm SS}|\).

\section{Fiducial definition and cutflow}

The fixed-order fiducial definition follows the central features of the
ATLAS trilepton particle-level region:
\begin{itemize}[leftmargin=1.8em]
\item exactly two electron-flavour leptons and one muon with total charge
  \(q_W=\pm1\);
\item the two same-sign leptons have \(\pt\geq20\) GeV; the remaining lepton
  has \(\pt\geq10\) GeV;
\item \(|\eta(e)|<2.47\) and \(|\eta(\mu)|<2.5\);
\item if an OSSF dielectron pair exists, \(m_{ee}>12\) GeV and
  \(|m_{ee}-m_Z|>10\) GeV; same-sign dielectron events pass these two
  conditions automatically;
\item
  \(|m_{3\ell}-m_Z|>10\) GeV;
\item at least one accepted jet and at least one accepted ideal b jet.
\end{itemize}
A parton-level isolation proxy requires every selected lepton and central
accepted jet to satisfy
\begin{equation}
 \dr(\ell,j)>
 \min\left(0.4,\;0.04+\frac{10~\mathrm{GeV}}{\pt(\ell)}\right).
\end{equation}
This is intentionally stated as a proxy: exact experimental overlap removal
acts sequentially on dressed leptons, particle jets, and b-hadron ghost tags,
none of which exists in a pure partonic NLO QCD event. There is no
missing-momentum cut in the fiducial definition.

The ten cutflow bin centres are:
\begin{center}
\begin{tabular}{cl}
\toprule
Bin & Accumulated requirement\\
\midrule
1 & valid visible \(2e1\mu+3\nu+b\bar b\) object content\\
2 & three-lepton \(\pt\) and \(\eta\) acceptance\\
3 & OSSF mass above 12 GeV when an OSSF pair exists\\
4 & OSSF \(Z\) veto when an OSSF pair exists\\
5 & trilepton \(Z\) veto\\
6 & at least one accepted jet\\
7 & at least one accepted b jet\\
8 & lepton--jet separation; this defines the fiducial sample\\
9 & two distinct accepted b jets; reconstruction input tier\\
10 & successful three-neutrino reconstruction\\
\bottomrule
\end{tabular}
\end{center}

\section{Three-neutrino reconstruction}

Only the three charged leptons, their charges, two accepted b jets, and
\(\vec p_{\mathrm T}^{\,\mathrm{miss}}\) enter the reconstruction. A trial
first assigns one of the two same-sign leptons to the associated W. The
other same-sign lepton belongs to the same-charge top, while the unique
opposite-sign lepton belongs to the other top. Both choices are tested. For
each assignment, a trial chooses transverse momenta for the positive- and
negative-top neutrinos; momentum conservation fixes
\[
 \vec q_{{\rm assoc}\,\mathrm T}=\vec p_{\mathrm T}^{\,\mathrm{miss}}
 -\vec q_{t^+\mathrm T}-\vec q_{t^-\mathrm T}.
\]
Each equation
\((p_\ell+p_\nu)^2=(m_W^{\rm ref})^2\) supplies up to two \(p_z^\nu\)
roots. A negative discriminant is replaced by its real part. All root
combinations, both same-sign origin assignments, and both b-jet assignments
are evaluated with
\begin{align}
 \chi^2_{\rm reco}={}&
 \sum_{q=t,\bar t}\left(\frac{m_q-m_t^{\rm ref}}{15~\mathrm{GeV}}\right)^2
 +\sum_{V=W_t^+,W_{\bar t}^-,W_{\rm assoc}}
 \left(\frac{m_V-m_W^{\rm ref}}{10~\mathrm{GeV}}\right)^2\\
 &+\left(\frac{m_t-m_{\bar t}}{30~\mathrm{GeV}}\right)^2
 +\left(\frac{\sum_i|p_{z,\nu_i}|}{1500~\mathrm{GeV}}\right)^2
 +\left(\frac{\sum_i\pt(\nu_i)}{2000~\mathrm{GeV}}\right)^2 .
\end{align}
The last two terms only resolve weakly constrained large-momentum solutions.

The four transverse degrees of freedom are initialized at equal sharing of
the measured missing momentum. A \(3^4\) coarse grid is followed by twelve
four-coordinate pattern-search iterations. The procedure is deterministic,
does not use truth matching, and returns all three W candidates, top,
antitop, \(\ttbar\), and \(\ttW\). It is a transparent fixed-order
reconstruction prescription rather than an emulation of a detector-specific
likelihood fit.

\section{Spin and polarization observables}

The top-pair spin basis is constructed in the \(\ttbar\) rest frame. With
\(\hat p\) the boosted positive-beam direction, \(\hat k\) the top direction,
and \(c=\hat p\cdot\hat k\),
\begin{equation}
 \hat r=\operatorname{sign}(c)
 \frac{\hat p-c\hat k}{\sqrt{1-c^2}},\qquad
 \hat n=\operatorname{sign}(c)
 \frac{\hat p\times\hat k}{\sqrt{1-c^2}} .
\end{equation}
A deterministic orthogonal completion covers forward exceptional points.
The \(e^+\) and \(e^-\) are boosted to their parent rest frames and define
six angles \(c_a^+\), \(c_b^-\), \(a,b\in\{k,r,n\}\). All nine products
\(q_{ab}=c_a^+c_b^-\) are recorded. For the convention
\begin{equation}
 \frac1\sigma\frac{d^2\sigma}{dc_a^+dc_b^-}
 =\frac14[1+B_a^+c_a^+ +B_b^-c_b^- -C_{ab}c_a^+c_b^-],
\end{equation}
\(B_a^\pm=3\langle c_a^\pm\rangle\) and
\(C_{ab}=-9\langle q_{ab}\rangle\). The histograms themselves are raw
moments and can be translated to another sign convention.

The associated-W suite adds the associated charged-lepton angle with the
positive beam in the W rest frame, its charge-signed counterpart, its
opening angle with the same-sign top-decay lepton, and
\[
 {\cal T}_{3\ell}=(\hat\ell_{t^+}\times\hat\ell_{t^-})
 \cdot\hat\ell_{\rm assoc} .
\]
It also records \(q_W{\cal T}_{3\ell}\) and, in the \(\ttW\) frame,
\(q_W[\cos(\ell_{\rm assoc},t)-
\cos(\ell_{\rm assoc},\bar t)]\). At truth level
\(\ell_{\rm assoc}=\mu\); the reconstructed tier uses the mass-fit
assignment.

\section{Complete histogram catalogue}

Histogram numbers are stable implementation labels. Energies and momenta
are in GeV; angles and rapidities are dimensionless.

\subsection{Rates and production: histograms 1--34}
\begin{longtable}{@{}L{1.3cm}L{5.1cm}L{9.1cm}@{}}
\toprule IDs & Histograms & Definition and purpose\\ \midrule \endhead
1 & inclusive rate & One-bin total valid \(2e1\mu\) rate for either allowed charge pattern.\\
2 & inclusive associated-W charge & \(q_W=-1,+1\); charge ratio and relative asymmetry input.\\
3 & cutflow & Ten cumulative stages defined above.\\
4 & fiducial associated-W charge & Charge-separated fiducial rate.\\
5--8 & top/antitop \(\pt,y\) & Truth top and antitop after QCD-radiation assignment.\\
9--15 & \(\ttbar\) mass, \(\pt\), \(y\), \(\Delta y\), \(\Delta|y|\), \(\ddphi\), \(\dr\) & Standard pair-production spectra and asymmetry inputs.\\
16--19 & scattering cosine, speed, \(\rho_{tt}\), \(\pt\) asymmetry & \(\cos\theta_t^*\), \(|\vec p_t^*|/E_t^*\), \(2m_t^{\rm ref}/m_{tt}\), and normalized top-\(\pt\) difference.\\
20--22 & associated-W \(\pt,y,m\) & Exact muon--neutrino W candidate.\\
23--25 & \(\ttW\) mass, \(\pt,y\) & Full heavy associated-production system.\\
26 & \(\ttW\) threshold \(\rho\) & \((2m_t^{\rm ref}+m_W^{\rm ref})/m_{ttW}\).\\
27 & heavy scalar \(\pt\) & \(\pt(t)+\pt(\bar t)+\pt(W_{\rm assoc})\).\\
28--31 & W--top, W--antitop, W--\(\ttbar\) \(\dr\); W--\(\ttbar\) \(\ddphi\) & Associated-W recoil geometry.\\
32--34 & \(q_Wy_W\), \(q_Wy_{ttW}\), \(q_W\Delta|y_t|\) & Charge-signed longitudinal and top-asymmetry projections.\\
\bottomrule
\end{longtable}

\subsection{Charged leptons: histograms 35--70}
\begin{longtable}{@{}L{1.3cm}L{5.1cm}L{9.1cm}@{}}
\toprule IDs & Histograms & Definition and purpose\\ \midrule \endhead
35--43 & \(e^+,e^-,\mu\) \(\pt,\eta,y\) & Individual charged leptons.\\
44--46 & ordered lepton \(\pt\) & Leading, subleading, and third lepton.\\
47--50 & trilepton \(m,\pt,y,H_T^\ell\) & Global visible leptonic system.\\
51--55 & OSSF \(ee\) \(m,\pt,\ddphi,|\Delta\eta|,\dr\) & Top-decay dilepton system and Z-veto variables.\\
56--60 & same-sign \(m,\pt,|\ddphi|,|\Delta\eta|,\dr\) & Signature pair; includes both ATLAS angular observables.\\
61--64 & opposite-sign \(e\mu\) \(m,\ddphi,|\Delta\eta|,\dr\) & Complementary pair geometry.\\
65--68 & min/max pair mass and min/max pair \(\dr\) & Permutation-insensitive three-lepton topology.\\
69--70 & \(q_W\eta_\mu,q_Wy_\mu\) & Charge-folded associated-lepton spectra.\\
\bottomrule
\end{longtable}

\subsection{W/top decays and neutrinos: histograms 71--110}
\begin{longtable}{@{}L{1.3cm}L{5.1cm}L{9.1cm}@{}}
\toprule IDs & Histograms & Definition and purpose\\ \midrule \endhead
71--79 & three W candidates, each \(\pt,y,m\) & Top-side, antitop-side, and associated W.\\
80--82 & three-W \(m,\pt,y\) & Sum of all three W candidates.\\
83--85 & correct \(e^+b\), \(e^-\bar b\), and maximum mass & Truth decay endpoints using bottom partons.\\
86--87 & electron energies in parent-top frames & Charged-lepton decay spectra.\\
88--90 & three W helicity cosines & Top, antitop, and associated-W polarization.\\
91 & \(q_W\cos\theta^*_{W,\rm assoc}\) & Common charge convention for W polarization.\\
92--93 & top-side W energies in parent frames & Top-decay dynamics.\\
94 & muon energy in associated-W frame & Off-shell-W and polarization sensitivity.\\
95--100 & three neutrino \(\pt,\eta\) & Exact truth neutrinos by decay chain.\\
101--103 & trineutrino \(\pt,m,y\) & Full invisible system.\\
104 & \(\met\) & Observable total missing transverse momentum.\\
105--106 & min/max \(\ddphi(\met,\ell)\) & Missing-momentum alignment.\\
107 & \(m_T(3\ell,\met)\) & Trilepton cluster transverse mass.\\
108--109 & \(m_T(\ell_{\rm lead},\met)\), \(m_T(\mu,\met)\) & Measurable transverse masses.\\
110 & missing-\(\pt\) closure & \(|\vec\met-\sum_\nu\vec\pt_\nu|\).\\
\bottomrule
\end{longtable}

\subsection{Jets and fiducial observables: histograms 111--178}
\begin{longtable}{@{}L{1.3cm}L{5.1cm}L{9.1cm}@{}}
\toprule IDs & Histograms & Definition and purpose\\ \midrule \endhead
111--116 & b/\(\bar b\) \(\pt\), ordered b \(\pt\), individual \(\eta\) & Ideal partonic b-tag kinematics.\\
117--120 & bb \(m,\pt,\ddphi,\dr\) & Two-b topology; merged bottoms have zero separation.\\
121--122 & \(N_j(\pt>25)\), \(N_j(\pt>40)\) & Central jet multiplicities.\\
123--126 & leading/subleading jet \(\pt,\eta\) & Ordered accepted jets.\\
127--130 & extra-jet multiplicity, leading-extra \(\pt,y\), \(H_T^{\rm extra}\) & Radiation beyond the b jets.\\
131--134 & \(H_T^j,H_T^\ell,H_T^{\rm vis},S_T\) & Event activity; first two are ATLAS unfolded variables.\\
135--136 & visible \(3\ell bb\) mass and \(\pt\) & Global reconstructible visible system.\\
137 & \(\dr_{\ell b,\rm lead}\) & Leading lepton to leading accepted b jet; ATLAS variable.\\
138 & minimum lepton--b \(\dr\) & Closest of six lepton--b combinations.\\
139--140 & fiducial rate and charge & Baseline accepted cross sections.\\
141--143 & fiducial ordered lepton \(\pt\) & Three accepted leptons.\\
144--146 & fiducial leading/trailing electron and muon \(\eta\) & Detector-acceptance shapes; electron ordering is by \(\pt\).\\
147--150 & fiducial trilepton \(m,\pt,y,H_T^\ell\) & Accepted trilepton system.\\
151--153 & fiducial dielectron \(m,\ddphi,\dr\) & Defined for opposite- and same-sign dielectron events; the Z veto acts only on OSSF pairs.\\
154--157 & fiducial same-sign \(m,|\ddphi|,|\Delta\eta|,\dr\) & Signature pair and ATLAS angles.\\
158--160 & fiducial leading opposite-sign pair \(m,\ddphi,\dr\) & The unique opposite-charge lepton paired with the leading same-sign lepton.\\
161--166 & \(\met\), min/max \(\ddphi\), leading/muon \(m_T\), cluster \(m_T\) & Neutrino-blind transverse observables.\\
167--170 & ordered b \(\pt\), bb mass and \(\dr\) & Visible heavy-flavour system.\\
171--175 & jet/extra-jet multiplicity, \(H_T^j,H_T^{\rm vis},S_T\) & Accepted activity.\\
176 & fiducial \(\dr_{\ell b,\rm lead}\) & Measurement-facing ATLAS variable.\\
177 & fiducial \(m_{\mathrm T2}^{ee}\) & Uses total \(\met\), including associated neutrino.\\
178 & accepted b-jet multiplicity & One-b fiducial versus two-b reconstruction tiers.\\
\bottomrule
\end{longtable}

\subsection{Reconstruction and spin: histograms 179--278}
\begin{longtable}{@{}L{1.3cm}L{5.1cm}L{9.1cm}@{}}
\toprule IDs & Histograms & Definition and purpose\\ \midrule \endhead
179 & reconstructed rate & Successful solutions within the two-b tier.\\
180--188 & reconstructed three W candidates, each \(\pt,y,m\) & Mass-constraint-fit W bosons.\\
189--194 & reconstructed top/antitop \(\pt,y,m\) & Both same-sign origins and both b-jet assignments tested.\\
195--200 & reconstructed \(\ttbar\), \(\ttW\), each \(m,\pt,y\) & Full fitted heavy systems.\\
201--206 & reconstructed neutrino \(\pt,p_z\) & Three fitted invisible components.\\
207 & reconstructed trineutrino mass & Fitted invisible-system invariant mass.\\
208--209 & fit score and b pairing & Minimum score; pairing 1 couples the positive/negative top leptons to b/\(\bar b\), pairing 2 swaps the jets.\\
210--212 & reconstructed W helicity cosines & Polarization after neutrino reconstruction.\\
213 & reconstructed \(q_W\cos\theta^*_{W,\rm assoc}\) & Charge-folded associated-W polarization.\\
214 & reconstructed \(\ttW\) threshold \(\rho\) & Fitted heavy-threshold shape.\\
215--220 & truth \(c_k^+,c_r^+,c_n^+,c_k^-,c_r^-,c_n^-\) & Complete single-spin angles.\\
221--229 & truth \(q_{ab}\), \(a,b\in\{k,r,n\}\) & All nine entries of the full \(3\times3\) correlation matrix.\\
230--240 & opening/\(\ddphi\), symmetric and antisymmetric moments, triple product, diagonal estimator, signed sine & Derived spin and CP-sensitive projections.\\
241--266 & reconstructed version of 215--240 & Same definitions using fitted tops.\\
267--272 & truth associated-spin suite & Beam cosines, same-sign opening, raw/signed triple product, and associated-lepton--top contrast.\\
273--278 & reconstructed associated-spin suite & Same six quantities in reconstructed frames.\\
\bottomrule
\end{longtable}

\subsection{Diagnostics: histograms 279--286}
\begin{longtable}{@{}L{1.3cm}L{5.1cm}L{9.1cm}@{}}
\toprule ID & Histogram & Meaning\\ \midrule \endhead
279 & truth top-assignment score & Minimum \(\chi^2_{\rm truth}\).\\
280 & assigned-radiation count & Light QCD partons assigned to either top decay.\\
281 & truth top-mass residual & Larger top/antitop reference-mass deviation.\\
282 & missing-\(\pt\) closure residual & Numerical event-record consistency.\\
283 & reconstruction score & Minimum \(\chi^2_{\rm reco}\), repeated for monitoring.\\
284 & reconstructed W-mass residual & Largest of three W reference-mass deviations.\\
285 & reconstructed top-mass residual & Larger fitted top-mass deviation.\\
286 & reconstruction rate & One-bin normalization for successful fits.\\
\bottomrule
\end{longtable}

\section{IR safety and fixed-order use}

The analysis never branches on \code{ibody}. The same code sees Born, real,
soft, collinear, and soft-collinear projected points in additive and
multiplicative decay-chain calculations. All measurement inputs are
infrared-safe functions of final momenta: anti-\(k_T\) jets, momentum sums,
ordered objects away from measure-zero ties, and ordinary acceptance cuts.
In particular, collinear \(\xi=0,y=1\) projections are experimentally
indistinguishable and are evaluated identically. Histogram boundaries,
jet-assignment boundaries, and mass-fit solution crossings move an event
only on measure-zero surfaces, as in any binned IR-safe analysis.

For a meaningful charge ratio, both \(W^+\) and \(W^-\) processes must be
included with identical settings. If only one charge is generated, the charge
histograms remain valid bookkeeping but do not define the hadronic asymmetry.

\section{Parton-level limitations}

This module is more realistic than a no-cut truth analysis but is not a
detector simulation. Leptons are bare QCD-level particles rather than
photon-dressed objects; b tags use bottom partons rather than ghost-associated
b hadrons; there is no detector efficiency, resolution, pile-up, fake lepton,
or charge-misidentification model. The \(2e1\mu\) flavour channel is one
component of an experimental inclusive \(e/\mu\) result. Reconstruction solutions are
prescription-dependent, and truth residuals are validation plots rather than
measurable observables.

The fiducial cuts are a reproducible theory region, not a claim of bitwise
identity with an ATLAS reconstruction. Direct data comparisons require a
matched particle-level definition, branching-ratio and flavour factors,
electroweak/QED treatment, and non-perturbative corrections.

\section{Implementation and selection}

Select the module with
\begin{verbatim}
FO_ANALYSIS_FORMAT=HwU
FO_ANALYSE=analysis_HwU_pp_ttxw_trilepton.o
\end{verbatim}
The fNLO makefile automatically adds
\code{analysis_HwU_pp_ttxw_trilepton_bridge.o}. The bridge exports the
ordinary \code{analysis_fill} ABI and variable-length
\code{analysis_fill_multiplicative} ABI. FastJet is linked through the
standard fNLO wrapper.

A representative \(W^+\) process is
\begin{verbatim}
p p > t t~ w+ [QCD],
 (t  > w+ b  QED=1 [QCD], w+ > e+ ve),
 (t~ > w- b~ QED=1 [QCD], w- > e- ve~),
 w+ > mu+ vm
\end{verbatim}
with the charge-conjugate associated-W decay for \(W^-\). Exact syntax can
depend on the model and decay-chain generation mode; object mapping in the
analysis does not depend on written order. The visible and reconstructed
histograms also accept, for example, a \(W^+\) configuration with
\(e^+e^+\mu^-\). The flavour-specific truth histograms remain empty for
such an alternative origin assignment, by construction.

\section{References}

\begin{enumerate}[leftmargin=1.8em]
\item ATLAS Collaboration, \emph{Measurement of the total and differential
  cross-sections of \(t\bar tW\) production in \(pp\) collisions at
  \(\sqrt{s}=13\) TeV with the ATLAS detector}, JHEP 05 (2024) 131,
  \href{https://arxiv.org/abs/2401.05299}{arXiv:2401.05299}.
\item G. Bevilacqua et al., \emph{NLO QCD corrections to off-shell
  \(t\bar tW^\pm\) production at the LHC: Correlations and Asymmetries},
  Eur. Phys. J. C 81 (2021) 675,
  \href{https://arxiv.org/abs/2012.01363}{arXiv:2012.01363}.
\item W. Bernreuther, D. Heisler and Z.-G. Si, \emph{A set of top quark
  spin correlation and polarization observables for the LHC},
  JHEP 12 (2015) 026,
  \href{https://arxiv.org/abs/1508.05271}{arXiv:1508.05271}.
\item M. Cacciari, G. P. Salam and G. Soyez, \emph{FastJet user manual},
  Eur. Phys. J. C 72 (2012) 1896,
  \href{https://arxiv.org/abs/1111.6097}{arXiv:1111.6097}.
\item M. Cacciari, G. P. Salam and G. Soyez, \emph{The anti-\(k_t\) jet
  clustering algorithm}, JHEP 04 (2008) 063,
  \href{https://arxiv.org/abs/0802.1189}{arXiv:0802.1189}.
\end{enumerate}

\end{document}
